\documentclass[a4paper,14pt]{extarticle}
\usepackage{suai-report}

\suaisetup{
  department   = 41,
  teacher-post = Старший преподаватель,
  teacher      = В. В. Боженко,
  type         = Отчет о лабораторной работе,
  number       = 3,
  title        = Регрессионный анализ,
  course       = Введение в анализ данных,
  group        = 4414,
  student      = Г. В. Моисеев,
  date         = today,
}

\begin{document}

\maketitle
\suaitoc

\suaiintro

Цель работы — изучение алгоритмов регрессии на практике.

\suaitasks
загрузить данные варианта, вывести описательные статистики и рассчитать парные коэффициенты корреляции Пирсона
обучить модели парной линейной регрессии по каждому из двух признаков, выполнить предсказание и рассчитать метрики качества
построить линии регрессии, график разницы предсказанных и истинных значений и график предсказаний относительно истинных значений
выбрать признак для полиномиальной регрессии, обучить модели нескольких степеней и сравнить их метрики на обучающей и тестовой выборках
загрузить набор данных salary.csv, выделить целевую переменную, обработать выбросы и закодировать категориальные признаки
построить матрицу диаграмм рассеяния, разделить данные на обучающую и тестовую выборки и масштабировать признаки
обучить линейную регрессию, оценить её качество и важность признаков по коэффициентам
реализовать регрессию методом k-ближайших соседей и деревом решений и сравнить результаты всех методов

\section{Вариант задания}

Номер в списке группы — 16. В файле regression.xlsx содержатся данные 10 вариантов, поэтому номер взят по кругу и выполнен вариант 6 (лист var6). Лист содержит 1000 наблюдений с двумя признаками x1 и x2 и целевой переменной y. Эти данные используются в частях 1 и 2.

В части 3 используется набор данных salary.csv: 1500 строк, каждая описывает одного сотрудника. Столбцы набора приведены в~\tabref{columns}.

\suaitable[columns]{Столбцы набора данных salary.csv}
Столбец | Содержание
age | возраст сотрудника, лет
income | годовой доход до текущей работы
experience | количество полных лет профессионального опыта
score\_test | результат профессионального теста
hours\_worked | среднее количество рабочих часов в неделю
distance\_to\_work | расстояние от дома до места работы
savings | сумма сбережений на банковском счёте
debt | общий объём текущих долгов
education | уровень образования: High School, Bachelor, Master, PhD
city | город проживания
job\_type | сфера профессиональной деятельности
marital\_status | семейное положение: Single, Married, Divorced
car | наличие автомобиля: Yes, No
remote\_work | работает ли сотрудник удалённо: Yes, No
salary | годовая зарплата

Работа выполнена в Jupyter-ноутбуке ЛР3\_Моисеев\_6.ipynb. Ссылка на ноутбук: \url{https://github.com/Echways/ItDA/tree/main/lab-3}.

\section{Простая линейная регрессия}

\subsection{Загрузка данных и корреляция}

Данные варианта загружены в DataFrame методом read\_excel библиотеки pandas. В таблице 1000 строк и 3 столбца типа float64, пропусков нет. Описательные статистики, полученные методом describe, показаны на~\figref{p1_describe}.

\suaiimg[0.43\textwidth]{p1_describe}{Описательные статистики данных варианта}

Признаки x1 и x2 лежат в пределах от $-2$ до 2, среднее и медиана равны нулю, квартили равны $-1$ и 1. Стандартное отклонение 1,155 совпадает со значением $4 / \sqrt{12}$ для равномерного распределения на отрезке длиной 4, то есть значения признаков распределены равномерно. Целевая переменная y меняется от $-5{,}67$ до 5,52, её среднее равно $-0{,}008$, стандартное отклонение 1,93. Масштаб признаков одинаков, поэтому стандартизация в частях 1 и 2 не выполнялась.

Парные коэффициенты корреляции Пирсона приведены на~\figref{p1_corr}.

\suaiimg[0.38\textwidth]{p1_corr}{Коэффициенты корреляции Пирсона}

Корреляция x1 и y равна 0,77: связь прямая и сильная. Корреляция x2 и y равна 0,51: связь прямая, средней силы. Признаки между собой не связаны, коэффициент равен $-0{,}003$. По этим значениям линейная модель по x1 должна быть точнее, чем по x2. При этом коэффициент Пирсона оценивает только линейную связь, поэтому значение 0,51 не исключает, что связь x2 и y нелинейна.

\subsection{Обучение моделей и предсказание}

Данные разделены на обучающую и тестовую выборки функцией train\_test\_split в соотношении 3 к 1: 750 и 250 строк. Разбиение одно для обеих моделей, поэтому модели проверяются на одних и тех же строках и их метрики сравнимы. Затем обучены две модели парной линейной регрессии $y = ax + b$: по x1 и по x2 (\figref{p1_fit}).

\suaiimg[0.9\textwidth]{p1_fit}{Разделение выборки и обучение моделей}

Предсказания обеих моделей на тестовой выборке собраны в один датафрейм вместе с истинными значениями (\figref{p1_predict}).

\suaiimg[0.83\textwidth]{p1_predict}{Истинные и предсказанные значения y}

Предсказания по x2 лежат ближе к нулю, чем истинные значения: при $y = -4{,}13$ модель по x1 дала $-2{,}36$, а модель по x2 — только $-1{,}65$.

\subsection{Метрики качества}

Качество моделей оценено четырьмя метриками. Средняя квадратичная ошибка MSE вычисляется по~\formref{mse}, средняя абсолютная ошибка MAE — по~\formref{mae}, коэффициент детерминации $R^2$ — по~\formref{r2}. RMSE равна корню из MSE и измеряется в тех же единицах, что и целевая переменная.

\suaieq[mse]{MSE = \frac{1}{n} \sum_{i=1}^{n} (y_i - \hat{y}_i)^2}
n | количество наблюдений
y_i | истинное значение целевой переменной
\hat{y}_i | предсказанное значение

\suaieq[mae]{MAE = \frac{1}{n} \sum_{i=1}^{n} |y_i - \hat{y}_i|}

\suaieq[r2]{R^2 = 1 - \frac{\sum_{i=1}^{n} (y_i - \hat{y}_i)^2}{\sum_{i=1}^{n} (y_i - \bar{y})^2}}
\bar{y} | среднее истинных значений

MSE сильно штрафует большие ошибки, MAE показывает, на сколько в среднем предсказание отличается от истинного значения. $R^2$ — доля дисперсии целевой переменной, объяснённая моделью: 1 у идеальной модели и 0 у модели, которая всегда предсказывает среднее. Метрики обеих моделей на тестовой выборке показаны на~\figref{p1_metrics}.

\suaiimg[\textwidth]{p1_metrics}{Метрики качества моделей парной регрессии}

Модель по x1 объясняет 56~\% дисперсии y, модель по x2 — только 20~\%. Средняя абсолютная ошибка модели по x1 равна 0,93, по x2 — 1,36, то есть в полтора раза больше. У обеих моделей RMSE больше MAE (1,22 против 0,93 и 1,64 против 1,36): в данных есть отдельные большие ошибки, которые возведение в квадрат усиливает. Стандартное отклонение y в тестовой выборке равно 1,83, и RMSE модели по x2 близка к этому значению, то есть такая модель немногим лучше предсказания средним. Ни одна из моделей не точна, так как y зависит от обоих признаков, а каждая модель использует только один.

\subsection{Коэффициенты регрессии}

Коэффициенты a и b обученных моделей приведены на~\figref{p1_coef}.

\suaiimg[0.84\textwidth]{p1_coef}{Коэффициенты a и b}

Получены уравнения $y = 1{,}308 x_1 + 0{,}027$ и $y = 0{,}881 x_2 - 0{,}024$. При увеличении x1 на единицу y в среднем растёт на 1,31, при увеличении x2 — на 0,88. Свободный член у обеих моделей близок к нулю: средние значения признаков и y равны нулю, и прямые проходят почти через начало координат.

\subsection{Визуализация регрессии}

Диаграммы рассеяния всех наблюдений с наложенными линиями регрессии показаны на~\figref{p1_line}.

\suaiimg[\textwidth]{p1_line}{Наблюдения и линии парной регрессии}

Облако точек x1 вытянуто вдоль прямой, его ширина примерно одинакова на всём отрезке. Зависимость линейная, прямая описывает её правильно, а разброс вокруг прямой вызван влиянием x2 и шумом. Облако x2 изогнуто: в середине, от $-1$ до 1, значение y почти не меняется, а у краёв меняется быстро. Слева точки уходят ниже прямой, справа — выше, прямая эту форму не повторяет.

Разница предсказанного и истинного значения в каждой точке тестовой выборки показана на~\figref{p1_diff}. Точка выше красной линии означает, что модель завысила значение, ниже — занизила.

\suaiimg[\textwidth]{p1_diff}{Разница предсказанного и истинного значения}

У модели по x1 точки разбросаны вокруг нуля без какой-либо формы, от $-3{,}3$ до 4,5, у 63~\% точек ошибка по модулю меньше 1. Ошибка не зависит от x1, значит прямая описывает связь полностью. У модели по x2 разброс шире, от $-4{,}2$ до 4,1, ошибка меньше 1 только у 42~\% точек. В ошибках видна волна: при x2 от $-2$ до $-1{,}5$ модель в среднем завышает y на 0,91, от $-1{,}5$ до $-0{,}5$ — занижает на 0,79, от 0,5 до 1,5 — снова завышает на 0,58. Зависимость ошибки от признака указывает на нелинейную связь, которую прямая не описывает.

График предсказанных значений относительно истинных показан на~\figref{p1_pred_true}. Красная линия соответствует истинным значениям, то есть идеальному совпадению прогноза с фактом, синие точки — предсказания.

\suaiimg[\textwidth]{p1_pred_true}{Предсказанные и истинные значения}

У модели по x1 точки идут вдоль красной линии, но облако наклонено слабее: предсказания лежат в пределах от $-2{,}6$ до 2,6, а истинные значения — от $-4{,}5$ до 4,6. Крайние значения модель не достигает: большие значения y она занижает, малые завышает. У модели по x2 облако почти горизонтальное, предсказания не выходят за пределы от $-1{,}8$ до 1,7, модель слабо реагирует на изменение y.

\subsection{Выводы по части 1}

Связь x1 и y линейная и сильная: парная регрессия по x1 даёт $R^2 = 0{,}56$ и MAE~= 0,93, её ошибки не зависят от признака. Связь x2 и y слабее и нелинейна: парная регрессия по x2 даёт $R^2 = 0{,}20$ и MAE~= 1,36, а в её ошибках видна волна. Для прогноза y по одному признаку лучше подходит x1. Для x2 прямая является слишком простой моделью.

\section{Полиномиальная регрессия}

\subsection{Выбор признака}

Полиномиальная регрессия применена к признаку x2. На диаграмме рассеяния (\figref{p1_line}) облако x2 изогнуто, а облако x1 вытянуто вдоль прямой. В ошибках линейной модели по x2 есть волна (\figref{p1_diff}), в ошибках модели по x1 её нет.

Для обучения и оценки моделей написаны две функции (\lstref{polynomial}). Функция fit\_polynomial собирает модель из двух шагов: преобразователь PolynomialFeatures добавляет к признаку его степени, затем на них обучается LinearRegression. Столбец из единиц отключён параметром include\_bias, так как свободный член линейная регрессия подбирает сама. Функция train\_test\_scores вычисляет $R^2$ и MAE на обучающей и тестовой выборках.

\suaicode[Python, polynomial]{Функции обучения и оценки полиномиальной регрессии}
    def fit_polynomial(degree, feature='x2'):
        model = make_pipeline(
            PolynomialFeatures(degree=degree, include_bias=False),
            LinearRegression())
        return model.fit(train[[feature]], train['y'])


    def train_test_scores(model, X_train, X_test, y_train, y_test):
        pred_train = model.predict(X_train)
        pred_test = model.predict(X_test)
        return {'R2 train': r2_score(y_train, pred_train),
                'R2 test': r2_score(y_test, pred_test),
                'MAE train': mean_absolute_error(y_train, pred_train),
                'MAE test': mean_absolute_error(y_test, pred_test)}

Для проверки выбора на каждом признаке обучены полиномы степени 1 и степени 3 (\figref{p2_feature_check}).

\suaiimg[0.62\textwidth]{p2_feature_check}{Метрики полиномов степени 1 и 3 для каждого признака}

Для x1 переход к степени 3 ничего не даёт: $R^2$ на тестовой выборке меняется с 0,558 до 0,561. Для x2 он растёт с 0,200 до 0,273, а MAE уменьшается с 1,358 до 1,315. Далее используется признак x2.

\subsection{Полином степени 2}

Обучена модель степени 2 и выполнено предсказание на тестовой выборке (\figref{p2_fit}). Метрики модели показаны на~\figref{p2_metrics}, линия регрессии — на~\figref{p2_degree2}.

\suaiimg[0.87\textwidth]{p2_fit}{Обучение полинома степени 2 и предсказание}

\suaiimg[0.97\textwidth]{p2_metrics}{Метрики полинома степени 2}

\suaiimg[0.7\textwidth]{p2_degree2}{Полиномиальная регрессия степени 2}

Полином степени 2 не лучше прямой: $R^2$ на тестовой выборке равен 0,193 против 0,200 у линейной модели, а линия на графике почти прямая. Парабола симметрична относительно вершины, а зависимость y от x2 слева уходит вниз, справа — вверх. Такую форму даёт нечётная степень.

\subsection{Другие степени полинома}

Те же шаги повторены для степеней 1, 3, 10 и 25. Степень 1 соответствует линейной регрессии из части 1 и нужна для сравнения. Метрики показаны на~\figref{p2_scores} и сведены в~\tabref{degrees}.

\suaiimg[\textwidth]{p2_scores}{Метрики полиномов разных степеней}

\suaitable[degrees]{Метрики полиномиальной регрессии y по x2}
Степень полинома | $R^2$ train | $R^2$ test | MAE train | MAE test
1 | 0,271 | 0,200 | 1,406 | 1,358
2 | 0,272 | 0,193 | 1,404 | 1,364
3 | 0,334 | 0,273 | 1,355 | 1,315
10 | 0,339 | 0,280 | 1,343 | 1,306
25 | 0,359 | 0,245 | 1,318 | 1,321

Предсказания моделей на тестовой выборке приведены на~\figref{p2_predictions}, линии регрессии — на~\figref{p2_degrees}.

\suaiimg[0.83\textwidth]{p2_predictions}{Истинные значения и предсказания полиномов разных степеней}

\suaiimg[\textwidth]{p2_degrees}{Полиномиальная регрессия степеней 1, 3, 10 и 25}

\subsection{Выводы по части 2}

Оптимальная степень полинома — 3. При переходе от степени 2 к степени 3 $R^2$ на тестовой выборке вырос с 0,193 до 0,273, MAE уменьшилась с 1,364 до 1,315. Кривая повторяет форму облака: плато в середине и быстрое изменение у краёв.

Если степень слишком низкая, модель недообучается. Полиномы степеней 1 и 2 одинаково плохи и на обучающей, и на тестовой выборке ($R^2$ около 0,27 и 0,20): им не хватает гибкости, чтобы повторить изгиб.

Если степень слишком высокая, модель переобучается. У полинома степени 25 лучший $R^2$ на обучающей выборке (0,359), но на тестовой он упал до 0,245, ниже, чем у степени 3. Кривая на графике колеблется, особенно у краёв отрезка: она подстраивается под шум в обучающих точках, а не под зависимость.

Степень 10 даёт на тестовой выборке $R^2 = 0{,}280$, на 0,006 больше степени 3. Такой прирост лежит в пределах случайного разброса, а коэффициентов у модели становится 10 вместо 3, поэтому усложнение модели не оправдано. Даже лучший полином объясняет только 27~\% дисперсии y: целевая переменная сильнее зависит от x1, а эта модель x1 не использует.

\section{Решение задачи регрессии различными методами}

\subsection{Загрузка данных}

Набор данных salary.csv загружен методом read\_csv. Результат метода info показан на~\figref{salary_info}.

\suaiimg[0.57\textwidth]{salary_info}{Общая информация о наборе данных salary.csv}

В таблице 1500 строк и 15 столбцов, пропусков и дубликатов нет. Числовых столбцов 9, категориальных 6. Описательные статистики числовых столбцов показаны на~\figref{salary_describe}.

\suaiimg[0.9\textwidth]{salary_describe}{Описательные статистики числовых столбцов}

У столбца salary среднее (98~573) заметно больше медианы (80~310), стандартное отклонение (111~787) больше среднего, а максимум 985~625 в десять раз превышает третий квартиль: в столбце есть аномально большие значения. У столбца income минимум отрицательный ($-826{,}4$), хотя доход отрицательным быть не может. Столбцы distance\_to\_work, savings и debt скошены вправо.

\subsection{Целевая переменная}

Целевая переменная — salary, годовая зарплата на текущей работе. Столбец income близок по смыслу, но содержит доход до текущей работы: он известен заранее и остаётся признаком. Гистограмма и диаграмма размаха целевой переменной показаны на~\figref{salary_before}.

\suaiimg[\textwidth]{salary_before}{Гистограмма и диаграмма размаха зарплаты до обработки}

Основная масса зарплат собрана в узкий пик около 80~000, правее тянется редкий хвост до миллиона. На диаграмме размаха ящик сжат в полоску, справа от него расположены десятки выбросов.

\subsection{Обработка выбросов}

Границы выбросов найдены по правилу полутора межквартильных размахов (\figref{outliers}).

\suaiimg[0.9\textwidth]{outliers}{Границы и количество выбросов зарплаты}

За границами 36~501 и 124~719 лежат 48 значений, но они различаются по природе. Два значения снизу (28~589 и 34~670) и одно сверху (128~998) продолжают основное распределение и вышли за границы незначительно. Остальные 45 значений начинаются с 312~155, между ними и основной массой значений нет. Для этих 45 строк зарплата сравнена с доходом до текущей работы (\figref{anomalies}).

\suaiimg[\textwidth]{anomalies}{Проверка аномальных зарплат}

У обычных строк зарплата в 1,6 раза больше прежнего дохода, у аномальных — в 14 раз, при этом остальные признаки у аномальных строк обычные. Это ошибки в данных, а не высокооплачиваемые сотрудники, поэтому 45 строк удалены. Три граничных значения оставлены как настоящие, хотя и редкие зарплаты. Также удалена одна строка с отрицательным доходом, так как восстановить его значение нельзя. Целевая переменная после обработки показана на~\figref{salary_after}.

\suaiimg[\textwidth]{salary_after}{Гистограмма и диаграмма размаха зарплаты после обработки}

Осталось 1454 строки из 1500. Зарплата распределена симметрично и близка к нормальному распределению: значения лежат в пределах от 28~589 до 128~998, среднее 80~099 почти совпадает с медианой 80~025, стандартное отклонение 15~295.

По тому же правилу проверены числовые признаки (\figref{feature_outliers}).

\suaiimg[0.56\textwidth]{feature_outliers}{Количество выбросов в числовых признаках}

Больше всего выбросов у distance\_to\_work (64), savings (34) и debt (14). Это длинные правые хвосты распределений, а не ошибки: большинство сотрудников живёт близко к работе и имеет небольшие сбережения. У income и score\_test около десяти выбросов с правдоподобными значениями. Эти строки оставлены.

\subsection{Кодирование признаков}

Признак education порядковый: уровни образования идут по возрастанию. Он закодирован числами от 0 до 3 с помощью OrdinalEncoder, порядок уровней задан явно, иначе кодировщик расставил бы их по алфавиту. Признаки city, job\_type и marital\_status номинальные, порядка между их значениями нет, поэтому к ним применено прямое кодирование функцией get\_dummies. Признаки car и remote\_work принимают два значения и тем же способом преобразованы в один столбец из нулей и единиц.

Параметр drop\_first удаляет по одному столбцу у каждого признака. Без этого сумма фиктивных столбцов одного признака всегда равна единице, что вместе со свободным членом даёт строгую линейную зависимость. Удалённые значения стали базовыми: Chicago, Education, Divorced и No. Код и первые три строки результата показаны на~\figref{encode}.

\suaiimg[0.85\textwidth]{encode}{Кодирование категориальных признаков}

После кодирования в таблице 22 столбца: 21 признак и целевая переменная.

\subsection{Матрица диаграмм рассеяния}

Матрица диаграмм рассеяния построена функцией pairplot по числовым признакам и зарплате после обработки выбросов (\figref{pairplot}). Показана нижняя половина матрицы, так как верхняя её зеркально повторяет.

\suaiimg[\textwidth]{pairplot}{Матрица диаграмм рассеяния}

Нижняя строка матрицы показывает зависимость зарплаты от каждого признака. Зарплата растёт вместе со стажем experience и прежним доходом income, обе зависимости линейные. С остальными числовыми признаками зарплата образует облака без выраженного направления. Признаки между собой не связаны: на остальных диаграммах облака точек также не имеют направления.

\subsection{Разделение выборки и масштабирование}

Данные разделены на обучающую и тестовую выборки в соотношении 3 к 1: 1090 и 364 строки (\figref{split}).

\suaiimg[0.88\textwidth]{split}{Разделение на обучающую и тестовую выборки}

Признаки измеряются в разных единицах: доход — в десятках тысяч, образование — от 0 до 3. Поэтому они приведены к среднему 0 и стандартному отклонению 1 с помощью StandardScaler (\figref{scale}). После кодирования все признаки числовые, и масштабированы все: тогда коэффициенты линейной регрессии сравнимы между собой, а в методе k-ближайших соседей признаки вносят равный вклад в расстояние. Параметры масштабирования вычислены только по обучающей выборке, тестовая преобразована уже готовыми параметрами. Иначе сведения о тестовой выборке попали бы в обучение, то есть произошла бы утечка данных.

\suaiimg[0.9\textwidth]{scale}{Масштабирование признаков}

\subsection{Линейная регрессия}

Для оценки моделей написана функция evaluate (\lstref{evaluate}). Она вычисляет на тестовой выборке метрики MSE, MAE, RMSE и $R^2$, а также среднюю относительную ошибку MAPE по~\formref{mape}. Для контроля переобучения добавлен $R^2$ на обучающей выборке.

\suaieq[mape]{MAPE = \frac{1}{n} \sum_{i=1}^{n} \left| \frac{y_i - \hat{y}_i}{y_i} \right| \cdot 100\,\%}

\suaicode[Python, evaluate]{Оценка модели и обучение линейной регрессии}
    def evaluate(model):
        predicted = model.predict(X_test_scaled)
        mape = mean_absolute_percentage_error(y_test, predicted)
        return {**regression_metrics(y_test, predicted),
                'MAPE, %': 100 * mape,
                'R2 train': model.score(X_train_scaled, y_train)}


    linear = LinearRegression().fit(X_train_scaled, y_train)

Метрики обученной линейной регрессии показаны на~\figref{lr_metrics}.

\suaiimg[0.64\textwidth]{lr_metrics}{Метрики линейной регрессии}

Модель объясняет 86~\% дисперсии зарплаты на тестовой выборке ($R^2 = 0{,}860$). MAE равна 4885: в среднем прогноз отличается от фактической зарплаты на 4,9 тысячи. RMSE равна 6067 при стандартном отклонении зарплаты в тестовой выборке 16~225, то есть ошибка в 2,7 раза меньше разброса самой зарплаты. MAPE равна 6,38~\%. $R^2$ на обучающей выборке равен 0,852, почти как на тестовой, переобучения нет.

Истинные и предсказанные значения зарплаты приведены на~\figref{lr_predictions}.

\suaiimg[0.98\textwidth]{lr_predictions}{Истинные и предсказанные значения зарплаты}

Коэффициенты модели, упорядоченные по убыванию модуля, приведены на~\figref{lr_coef}. Признаки стандартизованы, поэтому коэффициент показывает, на сколько меняется зарплата при увеличении признака на одно стандартное отклонение, и чем он больше по модулю, тем важнее признак.

\suaiimg[0.85\textwidth]{lr_coef}{Коэффициенты линейной регрессии}

Важнее всего стаж: увеличение на одно стандартное отклонение, то есть на 11,5 года, добавляет к зарплате 10~290. Второй по важности признак — прежний доход (7526), третий — уровень образования (5350). Заметно слабее влияют работа в сфере IT (1719) и результат теста (1190). Коэффициенты остальных 16 признаков меньше 700 по модулю: возраст, сбережения, долги, расстояние до работы, город, семейное положение, автомобиль и удалённая работа на зарплату почти не влияют. Свободный член 79~982 равен средней зарплате в обучающей выборке.

Фактические и предсказанные значения показаны на~\figref{lr_plot}.

\suaiimg[0.6\textwidth]{lr_plot}{Фактическая и предсказанная зарплата, линейная регрессия}

Точки лежат узкой полосой вдоль красной линии на всём диапазоне зарплат. Систематического сдвига нет: модель одинаково точно предсказывает и низкие, и высокие зарплаты.

\subsection{Метод k-ближайших соседей и дерево решений}

Метод k-ближайших соседей предсказывает зарплату как среднее по k самым похожим сотрудникам обучающей выборки. Число соседей оставлено по умолчанию, $k = 5$. Дерево решений последовательно делит обучающую выборку по порогам признаков и в каждом листе предсказывает среднюю зарплату. Глубина дерева не ограничена, это также значение по умолчанию. Параметр random\_state фиксирует выбор между равноценными разбиениями. Обучение обеих моделей показано на~\figref{knn_tree}.

\suaiimg[0.94\textwidth]{knn_tree}{Обучение моделей k-ближайших соседей и дерева решений}

\subsection{Сравнение методов}

Метрики трёх моделей собраны в один датафрейм (\figref{results}) и приведены в~\tabref{results}.

\suaiimg[\textwidth]{results}{Метрики всех обученных моделей}

\suaitable[results]{Метрики моделей на тестовой выборке}
Модель | MAE | RMSE | MAPE, \% | $R^2$ | $R^2$ train
Линейная регрессия | 4885 | 6067 | 6,38 | 0,860 | 0,852
Дерево решений | 8086 | 10~140 | 10,40 | 0,608 | 1,000
k-ближайших соседей, $k = 5$ | 9429 | 11~789 | 12,23 | 0,471 | 0,648

Предсказания трёх моделей приведены на~\figref{models_predictions}, графики фактических и предсказанных значений — на~\figref{models_plot}.

\suaiimg[0.69\textwidth]{models_predictions}{Истинные значения и предсказания трёх моделей}

\suaiimg[\textwidth]{models_plot}{Фактическая и предсказанная зарплата для трёх моделей}

У линейной регрессии точки лежат вдоль красной линии узкой полосой. У метода k-ближайших соседей облако наклонено слабее линии: предсказания стянуты к средней зарплате и лежат в пределах от 57~805 до 107~120 при фактических зарплатах от 44~219 до 122~930. Низкие зарплаты модель завышает, высокие занижает. Расстояние вычисляется по всем 21 признаку с равным весом, а на зарплату заметно влияют только три, поэтому ближайшие соседи часто похожи по неважным признакам. У дерева решений облако идёт вдоль красной линии, но оно заметно шире, чем у линейной регрессии. Глубина не ограничена, и дерево делило выборку, пока не воспроизвело точно каждую обучающую зарплату: $R^2$ на обучающей выборке равен 1,000, а на тестовой только 0,608. Дерево запомнило обучающие примеры вместе с шумом, то есть переобучилось. Одна точка лежит далеко внизу: сотруднику с зарплатой 67~282 предсказано 28~589. Это самая низкая зарплата в наборе, сотрудник попал в лист с этим обучающим примером.

Все предсказания показаны на одном графике на~\figref{models_sorted}. Сотрудники тестовой выборки упорядочены по фактической зарплате, чёрная линия соответствует фактическим значениям.

\suaiimg[\textwidth]{models_sorted}{Фактическая зарплата и предсказания трёх моделей}

Точки линейной регрессии держатся рядом с чёрной линией по всей её длине. Предсказания метода k-ближайших соседей образуют пологую полосу: слева она выше линии, справа ниже. Предсказания дерева следуют за линией, но разбросаны вокруг неё шире, чем у линейной регрессии.

\subsection{Выводы по части 3}

Целесообразнее использовать линейную регрессию. По средней абсолютной ошибке она почти вдвое точнее метода k-ближайших соседей и в 1,7 раза точнее дерева решений, не переобучается и легко интерпретируется: коэффициенты прямо показывают вклад каждого признака. Причина в самих данных: зарплата линейно зависит от стажа, прежнего дохода и уровня образования, и линейная модель описывает эту зависимость точно. Дерево решений без ограничения глубины переобучается: $R^2$ на обучающей выборке равен 1,000, на тестовой 0,608. Методу k-ближайших соседей мешают 18 признаков, которые на зарплату почти не влияют, но участвуют в расчёте расстояния наравне с важными.

\suaiconclusion

В ходе работы изучены алгоритмы регрессии на двух наборах данных. Первый — данные варианта 6 из файла regression.xlsx: 1000 наблюдений с признаками x1, x2 и целевой переменной y. Второй — набор salary.csv со сведениями о 1500 сотрудниках: возраст, стаж, прежний доход, образование, сфера деятельности и другие признаки, целевая переменная — годовая зарплата.

В первой части обучены две модели парной линейной регрессии. Связь y с x1 линейная (корреляция 0,77), модель $y = 1{,}308 x_1 + 0{,}027$ объясняет 56~\% дисперсии на тестовой выборке при MAE~= 0,93. Связь y с x2 нелинейна: прямая по x2 даёт $R^2 = 0{,}20$, а в её ошибках видна волна. Во второй части для x2 подобран полином степени 3, с которым $R^2$ на тестовой выборке вырос до 0,27. Полиномы степеней 1 и 2 недообучаются, а полином степени 25 переобучается: на обучающей выборке он лучший ($R^2 = 0{,}359$), на тестовой хуже степени 3 (0,245 против 0,273).

В третьей части выполнена предобработка набора salary.csv. Удалены 45 строк с аномальной зарплатой от 312 до 986 тысяч, которая в 14 раз превышала прежний доход, и одна строка с отрицательным доходом, осталось 1454 строки. Уровень образования закодирован порядковыми числами, остальные категориальные признаки — прямым кодированием с удалением базового столбца, получен 21 признак. Выборка разделена на обучающую (1090 строк) и тестовую (364 строки), признаки стандартизованы по параметрам обучающей выборки.

Лучший результат показала линейная регрессия: $R^2 = 0{,}860$, MAE~= 4885, средняя относительная ошибка 6,38~\%. Дерево решений дало $R^2 = 0{,}608$, метод k-ближайших соседей при $k = 5$ — 0,471. Зарплату определяют три признака: стаж, прежний доход и уровень образования, у них наибольшие по модулю коэффициенты линейной регрессии. Остальные 18 признаков на зарплату почти не влияют.

Сложность модели должна соответствовать форме зависимости в данных: линейную связь точнее всего описывает линейная модель, изогнутую — полином невысокой степени. Более гибкие модели на этих данных проигрывают, так как подстраиваются под шум или расходуют гибкость на неважные признаки.

\suaisources
Учебно-методические указания к выполнению лабораторной работы №3 по дисциплине «Введение в анализ данных».

\end{document}
